\documentclass[14pt]{extarticle}
\usepackage{fontspec}
\usepackage{polyglossia}
\usepackage{amsmath}
\usepackage{amssymb}
\usepackage{xcolor}
\usepackage{fancyhdr}
\usepackage{graphicx}
\usepackage{listings}
\usepackage[most]{tcolorbox}
% \usepackage{setspace}
\usepackage{pifont}
\usepackage{enumitem}
\usepackage{titlesec}
\usepackage[bottom]{footmisc}
\usepackage{titling}
\usepackage{minted}
\usepackage{etoolbox}
\usepackage{array}
\usepackage{longtable} % page-breakable tables (must load before polyglossia/bidi RTL setup)
\usepackage{extsizes}
\usepackage{float}
\usepackage{booktabs}
\usepackage{hyperref}
\usepackage[onehalfspacing]{setspace}
\newif\ifwhiteblack
\ifdefined\setwhiteblack
  \setwhiteblack
\fi

\ifwhiteblack
  \usepackage[a4paper, landscape, top=1.25cm, bottom=1.25cm, left=1.25cm, right=1.25cm, headheight=1.5cm, headsep=0.5cm]{geometry}
\else
  \usepackage[a4paper, top=2.7cm, bottom=2.7cm, left=2.5cm, right=2.5cm, headheight=1.5cm, headsep=0.5cm]{geometry}
\fi
\usepackage{tikz}
\usepackage{datetime2}

\usetikzlibrary{automata, positioning, arrows.meta, arrows}
\providecommand{\LTcaptype}{table}

\newfontfamily\emoji{Segoe UI Emoji}

\pagestyle{fancy}

\setmainlanguage[numerals=western]{arabic}
\setotherlanguage{english}
\newfontfamily\arabicfont[Script=Arabic]{Calibri}
\newfontfamily\arabicfonttt[Script=Arabic]{Courier New}



\ifwhiteblack
\lstset{
  language=[Sharp]C,
  numbers=left,
  stepnumber=1,
  numbersep=8pt,
  frame=single,
  basicstyle=\ttfamily\small,
  keywordstyle=\bfseries,
  stringstyle=\ttfamily,
  commentstyle=\itshape
}
\else
\lstset{
  language=[Sharp]C,
  numbers=left,
  stepnumber=1,
  numbersep=8pt,
  frame=single,
  basicstyle=\ttfamily\small,
  keywordstyle=\color{blue},
  stringstyle=\color{red},
  commentstyle=\color{green!50!black}
}
\fi

\newif\ifdetailed
\ifdefined\setdetailed
  \setdetailed
\fi

\newif\ifwithsols
\ifdefined\setwithsols
  \setwithsols
\fi

\ifwhiteblack
  \AtBeginDocument{
    \LARGE
    \usemintedstyle{bw}
  }
\fi

% Default minted look: compact code that wraps long lines.
% Not applied in black-and-white landscape mode, which already uses \LARGE.
\ifwhiteblack\else
  \setminted{fontsize=\small, breaklines=true}
\fi

% unified tcolorboxes for articles
\ifwhiteblack
\tcbset{colback=white, colframe=black, fonttitle=\bfseries, boxrule=1pt}
\newtcolorbox{boxDef}[1][]{title={{\emoji📘} تعريف\ifx\\#1\\\else ~#1\fi :}}
\newtcolorbox{boxExercise}[1][]{title={{\emoji🧩} تمرين\ifx\\#1\\\else ~#1\fi :}}
\newtcolorbox{boxExample}[1][]{title={{\emoji📝} مثال\ifx\\#1\\\else ~#1\fi :}}
\newtcolorbox{boxNote}[1][]{title={{\emoji✨} ملاحظة\ifx\\#1\\\else ~#1\fi :}}
\newtcolorbox{boxAttention}[1][]{title={{\emoji🔔} تنبيه\ifx\\#1\\\else ~#1\fi :}}
\newtcolorbox{boxWarning}[1][]{title={{\emoji⚡} ملاحظة هامة\ifx\\#1\\\else ~#1\fi :}}
\newtcolorbox{boxSolution}[1][]{title={{\emoji✅} حل\ifx\\#1\\\else ~#1\fi :}}
\newtcolorbox{boxSymbol}[1][]{title={{\emoji🔣} رمز\ifx\\#1\\\else ~#1\fi :}}
\newtcolorbox{boxHint}[1][]{title={{\emoji💡} تلميح\ifx\\#1\\\else ~#1\fi :}}
\else
\tcbset{colback=white, colframe=black, fonttitle=\bfseries, boxrule=0.8pt}
\newtcolorbox{boxDef}[1][]{colback=blue!5!white,colframe=blue!75!black,
  title={{\emoji📘} تعريف\ifx\\#1\\\else ~#1\fi :}}
\newtcolorbox{boxExercise}[1][]{colback=cyan!5!white,colframe=cyan!70!black,
  title={{\emoji🧩} تمرين\ifx\\#1\\\else ~#1\fi :}}
\newtcolorbox{boxExample}[1][]{colback=yellow!5!white,colframe=orange!90!black,
  title={{\emoji📝} مثال\ifx\\#1\\\else ~#1\fi :}}
\newtcolorbox{boxNote}[1][]{colback=gray!10!white,colframe=black,
  title={{\emoji✨} ملاحظة\ifx\\#1\\\else ~#1\fi :}}
\newtcolorbox{boxAttention}[1][]{colback=magenta!10!white,colframe=magenta!80!black,
  title={{\emoji🔔} تنبيه\ifx\\#1\\\else ~#1\fi :}}
\newtcolorbox{boxWarning}[1][]{colback=red!5!white,colframe=red!75!black,
  title={{\emoji⚡} ملاحظة هامة\ifx\\#1\\\else ~#1\fi :}}
\newtcolorbox{boxSolution}[1][]{colback=green!5!white,colframe=green!60!black,
  title={{\emoji✅} حل\ifx\\#1\\\else ~#1\fi :}}
\newtcolorbox{boxSymbol}[1][]{colback=purple!5!white,colframe=purple!70!black,
  title={{\emoji🔣} رمز\ifx\\#1\\\else ~#1\fi :}}
\newtcolorbox{boxHint}[1][]{colback=teal!5!white,colframe=teal!60!black,
  title={{\emoji💡} تلميح\ifx\\#1\\\else ~#1\fi :}}
\fi


\ifwhiteblack
\tcbset{simplecode/.style={ colback=white, colframe=black, boxrule=1pt, arc=2pt, left=4pt,right=4pt,top=4pt,bottom=4pt}}
\else
\tcbset{simplecode/.style={ colback=gray!5, colframe=black!50, boxrule=0.4pt, arc=2pt, left=4pt,right=4pt,top=4pt,bottom=4pt}}
\fi
% Compact code boxes: tighter line spacing (the preamble uses onehalfspacing) and small vertical gaps.
% Must come AFTER the simplecode definition above.
\tcbset{simplecode/.append style={fontupper=\setstretch{0.95}, top=4pt, bottom=4pt, before skip=6pt, after skip=6pt}}
\newenvironment{boxCode}{\begin{tcolorbox}[simplecode]}{\end{tcolorbox}}

% Code cell used inside tables (keeps code left-to-right inside RTL tables).
% Use inside a fixed-width p{...} column: \begin{pycell}\begin{minted}...\end{minted}\end{pycell}
\newtcolorbox{pycell}{frame empty,colback=white,boxsep=0pt,left=2pt,right=2pt,top=2pt,bottom=2pt,width=\linewidth}

\newcolumntype{C}[1]{>{\centering\arraybackslash}p{#1}}

\makeatletter
\DTMnewdatestyle{dotted}{%
  \renewcommand{\DTMdisplaydate}[4]{\two@digits{##3}.\two@digits{##2}.##1}%
  \renewcommand{\DTMDisplaydate}{\DTMdisplaydate}%
}
\makeatother
\DTMsetdatestyle{dotted}

% redefine spaces after titles
\makeatletter
\renewcommand{\@maketitle}{%
  \begin{center}
    {\huge \bfseries \@title \par}%
    \vskip 0.2em % space between title and author
    {\large \@author \par}%
    % \vskip 0.2em % space between author and date
    % {\normalsize \@date \par}%
    {\small \hfill \begin{english}\DTMtoday\end{english} \par}%
  \end{center}
}
\makeatother

\fancyhf{} % clear default
\fancypagestyle{plain}{
  \fancyhf{}
  \fancyhead[L]{مدرسة التسامح الشاملة}
  % \fancyhead[L]{\includegraphics[height=1cm]{../../../images/logoTasamoh.png}}
  \fancyhead[R]{الأستاذ محمود اغبارية}
  \fancyfoot[C]{\thepage}
}

\fancyhead[L]{مدرسة التسامح الشاملة}
\fancyhead[R]{الأستاذ محمود اغبارية}
\fancyfoot[C]{\thepage}

\setcounter{tocdepth}{3} % only section subsection and subsubsection in TOC

\newcommand{\importpdfpage}[2]{%
    \thispagestyle{empty}
    \newgeometry{left=0cm,right=0cm,top=0cm,bottom=0cm}
    \includegraphics[page=#2,width=0.95\paperwidth,height=0.95\paperheight,keepaspectratio]{#1}
    \restoregeometry
}

\newcommand{\insertFullImg}[1]{%
    \noindent
    \makebox[\textwidth][c]{
        \includegraphics[width=0.9\paperwidth,keepaspectratio]{#1}%
    }%
}%


\ifwithsols
\title{حل ورقة تمرن $6$ للصف العاشر $10$ - أسئلة شاملة في جملة الشرط \textenglish{if}}
\else
\title{ورقة تمرن $6$ للصف العاشر $10$ - أسئلة شاملة في جملة الشرط \textenglish{if}}
\fi

\begin{document}

\maketitle
\thispagestyle{fancy}

\begin{enumerate}[itemsep=2.2em]

% ==============================================================
% السؤال 4: حذف المنزلة الوسطى وفحص القسمة
% ==============================================================
\item
اكتب برنامجًا يستقبل من المستخدم عددًا صحيحًا موجبًا مكوّنًا من \textbf{$3$ منازل} (بين $100$ و $999$). \\
يقوم البرنامج بتكوين \textbf{عدد جديد مكوّن من منزلتين} ناتج عن حذف منزلة العشرات الوسطى من العدد الأصلي (فمثلاً من العدد $632$ يتكوّن العدد الجديد $62$)، ويطبع هذا العدد الجديد. \\
بعد ذلك، يفحص البرنامج العلاقة بين العدد الجديد ومنزلة العشرات التي حُذفت، ويطبع إحدى الرسائل التالية:
\begin{itemize}[nosep]
    \item إذا كانت منزلة العشرات المحذوفة تساوي $0$، لا يمكن القسمة عليها، فيطبع البرنامج: \textenglish{"Middle digit is zero"}.
    \item وإلا، إذا كان العدد الجديد يقبل القسمة على منزلة العشرات المحذوفة بدون باقٍ، يطبع البرنامج: \textenglish{"Divisible"}.
    \item وإلا، يطبع البرنامج: \textenglish{"Not divisible"}.
\end{itemize}

\textbf{مثال $1$:} إذا أدخل المستخدم العدد $633$، يحذف البرنامج منزلة العشرات، فيكوّن العدد الجديد $63$ ويطبعه، ثم يطبع \textenglish{Divisible} لأنّ $63$ يقبل القسمة على منزلة العشرات المحذوفة، وهي $3$، بدون باقٍ.

\textbf{مثال $2$:} إذا أدخل المستخدم العدد $408$، يحذف البرنامج منزلة العشرات، فيكوّن العدد الجديد $48$ ويطبعه، ثم يطبع \textenglish{Middle digit is zero} لأنّ منزلة العشرات المحذوفة تساوي $0$.

\ifwithsols
\begin{boxSolution}
\begin{english}
\begin{minted}{csharp}
Console.WriteLine("Enter a 3-digit number:");
int num = int.Parse(Console.ReadLine());

int hundreds = num / 100;
int tens = (num / 10) % 10;
int units = num % 10;

int newNum = hundreds * 10 + units;
Console.WriteLine(newNum);

if (tens == 0)
{
    Console.WriteLine("Middle digit is zero");
}
else if (newNum % tens == 0)
{
    Console.WriteLine("Divisible");
}
else
{
    Console.WriteLine("Not divisible");
}
\end{minted}
\end{english}
\end{boxSolution}
\fi
\filbreak

% ==============================================================
% السؤال 5: تعريف وفحص خاصية على عدد من 4 منازل - العدد المتكافئ
% ==============================================================
\item
نسمّي العدد الصحيح الموجب المكوّن من $4$ منازل (بين $1000$ و $9999$) \textbf{«عددًا متكافئًا»} إذا تحقّق فيه الشرطان التاليان معًا:
\begin{enumerate}[nosep]
    \item \textbf{مجموع الخانتين الخارجيتين} (أي خانة الآلاف وخانة الآحاد) يساوي \textbf{مجموع الخانتين الداخليتين} (أي خانة المئات وخانة العشرات).
    \item خانة الآلاف \textbf{لا تساوي} خانة الآحاد.
\end{enumerate}
على سبيل المثال:
\begin{itemize}[nosep]
    \item العدد $3849$ هو عدد متكافئ لأنّ $3 + 9 = 12$ وكذلك $8 + 4 = 12$، وخانة الآلاف لا تساوي الآحاد ($3 \neq 9$).
    \item العدد $4624$ هو عدد غير متكافئ؛ فرغم أنّ $4 + 4 = 8$ و $6 + 2 = 8$، إلا أنّ خانة الآلاف تساوي خانة الآحاد ($4 = 4$).
    \item العدد $2519$ هو عدد غير متكافئ لأنّ $2 + 9 = 11$ بينما $5 + 1 = 6$.
\end{itemize}
اكتب برنامجًا يستقبل من المستخدم عددًا صحيحًا موجبًا مكوّنًا من $4$ منازل، ويطبع \\
 \textenglish{"Equivalent"} إذا كان العدد متكافئًا، و \textenglish{"Not Equivalent"} وإلا.

\ifwithsols
\begin{boxSolution}
\begin{english}
\begin{minted}{csharp}
Console.WriteLine("Enter a 4-digit number:");
int num = int.Parse(Console.ReadLine());

int units = num % 10;
int tens = (num / 10) % 10;
int hundreds = (num / 100) % 10;
int thousands = num / 1000;

int outerSum = thousands + units;
int innerSum = hundreds + tens;

if (outerSum == innerSum && thousands != units)
{
    Console.WriteLine("Equivalent");
}
else
{
    Console.WriteLine("Not Equivalent");
}
\end{minted}
\end{english}
\end{boxSolution}
\fi
\filbreak

% ==============================================================
% السؤال 6: تعريف وفحص المتوالية الحسابية لثلاثة أعداد
% ==============================================================
\item
نسمّي ثلاثة أعداد صحيحة $(a, b, c)$ \textbf{«متوالية حسابية»} إذا كان الفرق بين العدد الثاني والأول مساويًا للفرق بين العدد الثالث والثاني (أي $b - a = c - b$). \\
اكتب برنامجًا يستقبل من المستخدم $3$ أعداد صحيحة $a$ ثم $b$ ثم $c$، ويصنّفها إلى إحدى الحالات التالية:
\begin{itemize}[nosep]
    \item إذا كانت الأعداد الثلاثة تشكل متوالية حسابية \textbf{متزايدة} (أي $b - a = c - b$ والفرق أكبر تمامًا من $0$، مثل $4, 7, 10$)، يطبع البرنامج: \textenglish{"Increasing Arithmetic"}.
    \item إذا كانت الأعداد الثلاثة تشكل متوالية حسابية \textbf{متناقصة} (أي $b - a = c - b$ والفرق أصغر تمامًا من $0$، مثل $20, 15, 10$)، يطبع البرنامج: \textenglish{"Decreasing Arithmetic"}.
    \item إذا كانت الأعداد الثلاثة تشكل متوالية حسابية \textbf{ثابتة} (أي الأعداد الثلاثة متساوية، مثل $5, 5, 5$)، يطبع البرنامج: \textenglish{"Constant Arithmetic"}.
    \item في جميع الحالات الأخرى (مثل $2, 5, 9$)، يطبع البرنامج: \textenglish{"Not Arithmetic"}.
\end{itemize}

\textbf{مثال $1$:} إذا أدخل المستخدم الأعداد $4$ و $7$ و $10$، يطبع البرنامج \textenglish{Increasing Arithmetic} لأنّ $7 - 4 = 3$ و $10 - 7 = 3$، والفرق $3 > 0$.

\textbf{مثال $2$:} إذا أدخل المستخدم الأعداد $12$ و $8$ و $4$، يطبع البرنامج \textenglish{Decreasing Arithmetic} لأنّ $8 - 12 = -4$ و $4 - 8 = -4$، والفرق $-4 < 0$.

\ifwithsols
\begin{boxSolution}
\begin{english}
\begin{minted}{csharp}
Console.WriteLine("Enter 3 integers:");
int a = int.Parse(Console.ReadLine());
int b = int.Parse(Console.ReadLine());
int c = int.Parse(Console.ReadLine());

int diff1 = b - a;
int diff2 = c - b;

if (diff1 == diff2)
{
    if (diff1 > 0)
    {
        Console.WriteLine("Increasing Arithmetic");
    }
    else if (diff1 < 0)
    {
        Console.WriteLine("Decreasing Arithmetic");
    }
    else
    {
        Console.WriteLine("Constant Arithmetic");
    }
}
else
{
    Console.WriteLine("Not Arithmetic");
}
\end{minted}
\end{english}
\end{boxSolution}
\fi
\filbreak

% ==============================================================
% السؤال 7: حافلات الرحلة المدرسية (قسمة وباقي مع شروط)
% ==============================================================
\item
تنظّم المدرسة رحلة تعليمية، وتريد حجز حافلات لنقل الطلاب. تتسع كل \textbf{حافلة كبيرة} لـ $45$ طالبًا بالضبط. \\
اكتب برنامجًا يستقبل من المستخدم العدد الكلي للطلاب المسجلين للرحلة (عدد صحيح موجب). \\
يطبع البرنامج أولاً عدد \textbf{الحافلات الكبيرة الممتلئة بالكامل}، ثم عدد \textbf{الطلاب المتبقين} الذين لا يكفون لملء حافلة كبيرة كاملة. \\
بعد ذلك، يطبع البرنامج توصية بخصوص الطلاب المتبقين حسب القواعد التالية:
\begin{itemize}[nosep]
    \item إذا لم يتبقَّ أي طالب (الباقي يساوي $0$)، يطبع: \textenglish{"No extra transport needed"}.
    \item إذا كان عدد الطلاب المتبقين بين $1$ و $15$ طالبًا (شاملاً $1$ و $15$)، يطبع: \textenglish{"Need 1 minibus"}.
    \item إذا كان عدد الطلاب المتبقين أكثر من $15$ طالبًا، يطبع: \textenglish{"Need 1 extra large bus"}.
\end{itemize}

\textbf{مثال $1$:} إذا أدخل المستخدم عدد الطلاب $100$، يطبع البرنامج أنّ عدد الحافلات الكبيرة الممتلئة هو $2$، وأنّ عدد الطلاب المتبقين هو $10$، ثم يطبع \textenglish{Need 1 minibus}.

\textbf{مثال $2$:} إذا أدخل المستخدم عدد الطلاب $115$، يطبع البرنامج أنّ عدد الحافلات الكبيرة الممتلئة هو $2$، وأنّ عدد الطلاب المتبقين هو $25$، ثم يطبع \textenglish{Need 1 extra large bus}.

\ifwithsols
\begin{boxSolution}
\begin{english}
\begin{minted}{csharp}
Console.WriteLine("Enter total students:");
int students = int.Parse(Console.ReadLine());

int fullBuses = students / 45;
int remaining = students % 45;

Console.WriteLine("Full large buses: " + fullBuses);
Console.WriteLine("Remaining students: " + remaining);

if (remaining == 0)
{
    Console.WriteLine("No extra transport needed");
}
else if (remaining <= 15)
{
    Console.WriteLine("Need 1 minibus");
}
else
{
    Console.WriteLine("Need 1 extra large bus");
}
\end{minted}
\end{english}
\end{boxSolution}
\fi
\filbreak

% ==============================================================
% السؤال 8: مجال تعريف تعبير رياضي وحساب قيمته
% ==============================================================
\item
معطى التعبير الرياضي التالي:
\[ f(x) = \frac{x^3 - 8}{\sqrt{25 - x^2}} \]
اكتب برنامجًا يستقبل من المستخدم عددًا عشريًّا يمثّل قيمة $x$، ويفحص مجال تعريف التعبير قبل حسابه:
\begin{itemize}[nosep]
    \item إذا كان $x = 5$ أو $x = -5$ (حينها يصبح المقام صفرًا)، يطبع البرنامج: \\
    \textenglish{"ERROR: division by 0"}.
    \item إذا كان $x > 5$ أو $x < -5$ (حينها يصبح ما تحت الجذر سالبًا)، يطبع البرنامج: \\
    \textenglish{"ERROR: negative square root"}.
    \item في باقي الحالات (أي عندما $-5 < x < 5$)، يحسب البرنامج قيمة $f(x)$ ويطبعها.
\end{itemize}

\textbf{مثال $1$:} إذا أدخل المستخدم $x = -5$، يطبع البرنامج \textenglish{ERROR: division by 0} لأنّ المقام يساوي صفرًا.

\textbf{مثال $2$:} إذا أدخل المستخدم $x = 3$، يحسب البرنامج القيمة $f(x) = 4.75$ ويطبعها، لأنّ $\frac{3^3 - 8}{\sqrt{25 - 3^2}} = \frac{19}{\sqrt{16}} = \frac{19}{4} = 4.75$.

\ifwithsols
\begin{boxSolution}
\begin{english}
\begin{minted}{csharp}
Console.WriteLine("Enter x:");
double x = double.Parse(Console.ReadLine());

if (x == 5 || x == -5)
{
    Console.WriteLine("ERROR: division by 0");
}
else if (x > 5 || x < -5)
{
    Console.WriteLine("ERROR: negative square root");
}
else
{
    double numerator = Math.Pow(x, 3) - 8;
    double denominator = Math.Sqrt(25 - Math.Pow(x, 2));
    double fx = numerator / denominator;
    Console.WriteLine("f(x) = " + fx);
}
\end{minted}
\end{english}
\end{boxSolution}
\fi
\filbreak

% ==============================================================
% السؤال 9: حل المعادلة التربيعية حسب المميز دلتا
% ==============================================================
\item
عند حل معادلة تربيعية من الصورة $ax^2 + bx + c = 0$ (حيث $a \neq 0$)، نرمز للمميّز بـ $\Delta = b^2 - 4ac$ وعندها:
\begin{itemize}[nosep]
    \item إذا كان $\Delta < 0$: لا توجد حلول حقيقية للمعادلة، فيطبع البرنامج \textenglish{"No real solutions"}.
    \item إذا كان $\Delta = 0$: يوجد حلّ حقيقي واحد للمعادلة وهو:
    \( x = \frac{-b}{2a} \)
    \item إذا كان $\Delta > 0$: يوجد حلّان حقيقيان للمعادلة، وهما:
    \[ x_1 = \frac{-b + \sqrt{\Delta}}{2a} \qquad x_2 = \frac{-b - \sqrt{\Delta}}{2a} \]
\end{itemize}
اكتب برنامجًا يستقبل معاملات المعادلة التربيعية $a, b, c$ ويحسب حلول المعادلة ويطبعها مع رسالة مناسبة.

\textbf{مثال $1$:} عند استقبال المعاملات $1, -5, 6$ يطبع البرنامج:
\textenglish{"Two solutions: x1 = 3, x2 = 2"}.

\textbf{مثال $2$:} عند استقبال المعاملات $1, 2, 1$ يطبع البرنامج:
\textenglish{"One solution: x = -1"}.

\textbf{مثال $3$:} عند استقبال المعاملات $1, 0, 1$ يطبع البرنامج:
\textenglish{"No real solutions"}.

\ifwithsols
\begin{boxSolution}
\begin{english}
\begin{minted}{csharp}
Console.WriteLine("Enter a:");
double a = double.Parse(Console.ReadLine());
Console.WriteLine("Enter b:");
double b = double.Parse(Console.ReadLine());
Console.WriteLine("Enter c:");
double c = double.Parse(Console.ReadLine());

double delta = Math.Pow(b, 2) - 4 * a * c;
if (delta < 0)
{
    Console.WriteLine("No real solutions");
}
else if (delta == 0)
{
    double x = -b / (2 * a);
    Console.WriteLine("One solution: x = " + x);
}
else
{
    double x1 = (-b + Math.Sqrt(delta)) / (2 * a);
    double x2 = (-b - Math.Sqrt(delta)) / (2 * a);
    Console.WriteLine("Two solutions: x1 = " + x1 + ", x2 = " + x2);
}
\end{minted}
\end{english}
\end{boxSolution}
\fi
\filbreak

% ==============================================================
% السؤال 10: حساب فاتورة الكهرباء التراكمية حسب الشرائح
% ==============================================================
\item
اكتب برنامجًا يستقبل عدد الكيلوواط المستهلكة (عددًا صحيحًا موجبًا)، ويحسب التكلفة التراكمية حسب التعرفة التدريجية التالية:
\begin{itemize}[nosep]
    \item أول $100$ كيلوواط (من $1$ إلى $100$): الكلفة $0.5$ شيكل لكل كيلوواط.
    \item كل كيلوواط إضافي من $101$ إلى $200$ كيلوواط: الكلفة $0.7$ شيكل لكل كيلوواط.
    \item كل كيلوواط إضافي أكثر من $200$: الكلفة $1.0$ شيكل لكل كيلوواط.
\end{itemize}
على البرنامج أن يطبع المبلغ الإجمالي المستحق للدفع.

\textbf{مثال $1$:} إذا كان المدخل $150$ كيلوواط، فالمبلغ الإجمالي هو:
$100 \times 0.5 + 50 \times 0.7 = 85$.

\textbf{مثال $2$:} إذا كان المدخل $780$ كيلوواط، فالمبلغ الإجمالي هو:
$100 \times 0.5 + 100 \times 0.7 + 580 \times 1 = 700$.

\ifwithsols
\begin{boxSolution}
\begin{english}
\begin{minted}{csharp}
Console.WriteLine("Enter kWh:");
int kwh = int.Parse(Console.ReadLine());

double cost;
if (kwh <= 100)
{
    cost = kwh * 0.5;
}
else if (kwh <= 200)
{
    cost = 100 * 0.5 + (kwh - 100) * 0.7;
}
else
{
    cost = 100 * 0.5 + 100 * 0.7 + (kwh - 200) * 1.0;
}
Console.WriteLine("Total cost = " + cost);
\end{minted}
\end{english}
\end{boxSolution}
\fi
\filbreak

% ==============================================================
% السؤال 11: تكلفة شحن الطرود البريدية حسب المنطقة والوزن
% ==============================================================
\item
تحسب شركة بريد تكلفة شحن الطرود بناءً على \textbf{منطقة الشحن} و\textbf{وزن الطرد} بالكيلوغرام. \\
اكتب برنامجًا يستقبل من المستخدم حرفًا واحدًا (\textenglish{char}) يمثّل منطقة الشحن (\textenglish{'L'} للشحن المحلي، أو \textenglish{'I'} للشحن الدولي)، ثم عددًا عشريًا موجبًا (\textenglish{double}) يمثّل وزن الطرد بالكيلوغرام، ويحسب ويطبع تكلفة الشحن حسب القواعد التالية:
\begin{itemize}
    \item \textbf{في حالة الشحن المحلي (\textenglish{'L'}):}
    \begin{itemize}[nosep]
        \item إذا كان الوزن $5$ كغم أو أقل، تكون التكلفة الثابتة $25$ شيكلاً.
        \item إذا زاد الوزن عن $5$ كغم، يدفع الزبون $25$ شيكلاً عن أوّل $5$ كغم، بالإضافة إلى $4$ شواكل عن كل كيلوغرام إضافي فوق الـ $5$.
    \end{itemize}
    \item \textbf{في حالة الشحن الدولي (\textenglish{'I'}):}
    \begin{itemize}[nosep]
        \item إذا كان الوزن $5$ كغم أو أقل، تكون التكلفة الثابتة $80$ شيكلاً.
        \item إذا زاد الوزن عن $5$ كغم، يدفع الزبون $80$ شيكلاً عن أوّل $5$ كغم، بالإضافة إلى $15$ شيكلاً عن كل كيلوغرام إضافي فوق الـ $5$.
    \end{itemize}
\end{itemize}

\textbf{مثال $1$:} إذا أدخل المستخدم المنطقة الدولية \textenglish{'I'} والوزن $8$ كغم، يطبع البرنامج تكلفة الشحن $125$ شيكلاً، لأنّ $80 + (8 - 5) \times 15 = 125$.

\textbf{مثال $2$:} إذا أدخل المستخدم المنطقة المحلية \textenglish{'L'} والوزن $3.5$ كغم، يطبع البرنامج تكلفة الشحن $25$ شيكلاً، لأنّ الوزن لا يتجاوز $5$ كغم.

\ifwithsols
\begin{boxSolution}
\begin{english}
\begin{minted}{csharp}
Console.WriteLine("Enter zone (L or I):");
char zone = char.Parse(Console.ReadLine());
Console.WriteLine("Enter weight in kg:");
double weight = double.Parse(Console.ReadLine());

double cost;
if (zone == 'L')
{
    if (weight <= 5)
    {
        cost = 25;
    }
    else
    {
        cost = 25 + (weight - 5) * 4;
    }
}
else
{
    if (weight <= 5)
    {
        cost = 80;
    }
    else
    {
        cost = 80 + (weight - 5) * 15;
    }
}
Console.WriteLine("Shipping cost: " + cost);
\end{minted}
\end{english}
\end{boxSolution}
\fi
\filbreak

% ==============================================================
% السؤال 12: حساب طلبية مطعم مع تخفيض الكمية ورسوم التوصيل
% ==============================================================
\item
في مطعم يبيعون نوعين من الوجبات:
\begin{itemize}[nosep]
    \item وجبة عادية (يرمز لها بالحرف \textenglish{'R'}): سعر الوجبة الواحدة $25$ شيكلاً.
    \item وجبة خاصة (يرمز لها بالحرف \textenglish{'S'}): سعر الوجبة الواحدة $35$ شيكلاً.
\end{itemize}
تُحسب التكلفة النهائية للطلبية وفق القواعد المتسلسلة التالية:
\begin{enumerate}[nosep]
    \item إذا شملت الطلبية \textbf{$5$ وجبات أو أكثر}، يحصل الزبون على تخفيض قدره $20$ شيكلاً من ثمن الوجبات.
    \item إذا طلب الزبون خدمة التوصيل للمنزل (\textenglish{needDelivery} من نمط \textenglish{bool})، تُضاف رسوم توصيل قدرها $15$ شيكلاً، \textbf{إلا} إذا كان ثمن الطلبية (بعد احتساب التخفيض إن وُجد) $150$ شيكلاً أو أكثر، فيكون التوصيل مجانيًا.
\end{enumerate}
اكتب برنامجًا يستقبل نوع الوجبة (\textenglish{'R'} أو \textenglish{'S'})، وعدد الوجبات، وهل الطلبية تشمل توصيلاً (\textenglish{true/false})، ويطبع المبلغ النهائي للدفع.

\textbf{مثال $1$:} إذا كانت الوجبة من النمط \textenglish{'R'}، وعددها $6$، والتوصيل مطلوبًا، يطبع البرنامج المبلغ النهائي $145$ شيكلاً؛ فثمن الوجبات $150$، وبعد التخفيض يصبح $130$، ثم تضاف رسوم التوصيل $15$ شيكلاً.

\textbf{مثال $2$:} إذا كانت الوجبة من النمط \textenglish{'S'}، وعددها $5$، والتوصيل مطلوبًا، يطبع البرنامج المبلغ النهائي $155$ شيكلاً؛ فثمن الوجبات $175$، وبعد التخفيض يصبح $155$، ولذلك يكون التوصيل مجانيًا.

\ifwithsols
\begin{boxSolution}
\begin{english}
\begin{minted}{csharp}
Console.WriteLine("Enter meal type (R/S):");
char mealType = char.Parse(Console.ReadLine());
Console.WriteLine("Enter number of meals:");
int count = int.Parse(Console.ReadLine());
Console.WriteLine("Need delivery? (true/false):");
bool needDelivery = bool.Parse(Console.ReadLine());

int total;
if (mealType == 'R')
{
    total = count * 25;
}
else
{
    total = count * 35;
}

if (count >= 5)
{
    total = total - 20;
}

if (needDelivery && total < 150)
{
    total = total + 15;
}

Console.WriteLine("Total price: " + total);
\end{minted}
\end{english}
\end{boxSolution}
\fi
\filbreak

\end{enumerate}

\vspace{1cm}
\begin{flushleft}
أرجو لكم وقتًا ممتعًا.

الأستاذ محمود اغبارية.
\end{flushleft}

\end{document}
